%%% Intro.tex --- 
%% 
%% Filename: Intro.tex
%% Description: 
%% Author: Ola Leifler
%% Maintainer: 
%% Created: Thu Oct 14 12:54:47 2010 (CEST)
%% Version: $Id$
%% Version: 
%% Last-Updated: Thu May 19 14:12:31 2016 (+0200)
%%           By: Ola Leifler
%%     Update #: 5
%% URL: 
%% Keywords: 
%% Compatibility: 
%% 
%%%%%%%%%%%%%%%%%%%%%%%%%%%%%%%%%%%%%%%%%%%%%%%%%%%%%%%%%%%%%%%%%%%%%%
%% 
%%% Commentary: 
%% 
%% 
%% 
%%%%%%%%%%%%%%%%%%%%%%%%%%%%%%%%%%%%%%%%%%%%%%%%%%%%%%%%%%%%%%%%%%%%%%
%% 
%%% Change log:
%% 
%% 
%% RCS $Log$
%%%%%%%%%%%%%%%%%%%%%%%%%%%%%%%%%%%%%%%%%%%%%%%%%%%%%%%%%%%%%%%%%%%%%%
%% 
%%% Code:



\chapter{Introduction}
\label{cha:introduction}
Modern data infrastructures, supercomputer systems and distributed systems generate massive amounts of real-time log data. These logs contain valuable insights into system health, performance, and potential security threats. However, manually analyzing such large and complex datasets is impractical due to their scale. As systems become increasingly interconnected and data volumes grow exponentially, the need for efficient and automated methods to process and analyze log data becomes critical.

Anomaly detection plays a key role in identifying irregular patterns in log data, which may indicate system failures, cyberattacks, or performance degradation. While anomaly detection is widely applied to network traffic analysis, its use in log data remains comparatively underexplored. Using anomaly detection on log data can fill two main purposes, looking back on offline log files to detect previous anomalies, and detecting anomalies on online log data in real-time. Online log data must be processed dynamically, requiring methods that can operate in real-time without relying on complete historical data. In contrast, offline log data can be processed statically, allowing for deeper analysis with access to the entire dataset.

 Machine learning based anomaly detection can enhance anomaly detection specifically on log data, reducing the workload. Also, machine learning in online streaming data enables an automated evaluation of logs, noticing threats before it is to late. This task comes with challenges. As online log data is generated at high speed, the anomaly detector requires to be fast enough to process and analyze data in real-time. Additionally, anomalies are not always obvious, there exist rare anomalies which can go unnoticed when they are compared to normal data. These are difficult for machine learning techniques to detect as they are normally constructed in a generalized way. This paper will analyze different anomaly detection methods within both offline and online log data, comparing the results between the two cases, as well as within each case. 

\section{Motivation}
\label{sec:motivation}
While numerous anomaly detection methods exist and have been widely studied, few studies compare the performance of anomaly detection techniques on both offline and online log data. Existing literature typically focuses on either historical log data or streaming log data, but rarely examines the combination of both. Additionally, the variation in performance for the five algorithms is rarely examined when applied on different types of datasets. Since log data can be utilized for both real-time and past-time anomaly detection, this project aims to evaluate both approaches to gain a deeper understanding of their differences.

\section{Aim}
\label{sec:aim}
In this thesis project, we analyze five anomaly detection techniques on different non-encrypted log datasets. Our objective is to assess and compare their effectiveness in detecting anomalies in different types of data. We focus on two distinct use cases where one involves analyzing historical log data offline, and the other involves processing streamed log data online. 

\section{Research questions}
\label{sec:research-questions}
At a high level, this thesis aims to identify the most effective anomaly detection technique for the given dataset and use case in form of online or offline analysis. To achieve this, five algorithms are tested on two log datasets, allowing for answers to the following questions:

\begin{enumerate}
\item What anomaly detection techniques are the most effective for identifying anomalies offline in different types of non-encrypted log datasets?

\item What anomaly detection techniques are the most effective for identifying anomalies online in different types of non-encrypted log datasets?

\item How does the anomaly detection results on online and offline log data differ, and why? 

\end{enumerate}

The answers to these questions will provide insights into selecting the most suitable anomaly detection techniques for the two use cases, online and offline. Likewise, it will provide information on what algorithm to use for different datasets.

\section{Approach}
\label{sec:approach}
To answer the research questions, this study uses log datasets, which hold suitability for both offline and online streaming data. An iterative approach is applied, refining the methods through multiple attempts to achieve the best results. This study will focus on machine learning-based anomaly detection techniques, evaluating the results of these techniques. The same framework is used for offline and online anomaly detection evaluation. This framework first converts the raw log file into structured logs, then extracts relevant features, which are used for anomaly detection with selected algorithms. For online log data however, this process is done for each log row before handling the next. In contrast, offline learning handles the entire file in each step. The results are evaluated based on recall, precision and F1-score.

\section{Contributions}
\label{sec:contributions}

\begin{itemize}
\item A systematic evaluation of five anomaly detection methods in non-encrypted log datasets to determine their effectiveness for offline and online log data.
\item A comparison and evaluation of how online and offline data impact the results of anomaly detection. We analyze the differences in performance and accuracy, when applying anomaly detection to streaming log data in real-time compared to processing historical log files offline. 
\item Analyzing anomaly detection techniques based on different key values to provide a deeper understanding of the strengths and weaknesses of each technique, allowing informed selection based on the specific use case and purpose.
\end{itemize}

\section{Thesis outline}
To analyze anomaly detection in log data, it is essential to first establish the necessary background, which is provided in Chapter \ref{cha:background}. Section \ref{logdata} introduces log data, followed by an overview of the datasets used in this study in Section \ref{Datasets}. Since the techniques used are based on machine learning, fundamental concepts in this area are covered in Section \ref{ML}, leading to an introduction to anomaly detection in Section \ref{AD}. We then outline earlier work and research on the subject in Section \ref{RelatedWork}. The process of parsing, feature extracting, and anomaly detecting the data is explored in Chapter \ref{cha:method}. The findings are then presented in Chapter \ref{cha:evaluation} and discussed in Chapter \ref{cha:discussion}. Finally, we conclude in Chapter \ref{cha:conclusion}.

%\nocite{scigen}
%We have included Paper \ref{art:scigen}

%%%%%%%%%%%%%%%%%%%%%%%%%%%%%%%%%%%%%%%%%%%%%%%%%%%%%%%%%%%%%%%%%%%%%%
%%% Intro.tex ends here


%%% Local Variables: 
%%% mode: latex
%%% TeX-master: "demothesis"
%%% End: 
